\scanpage{63}

\mspage{4}

relations (1) est plat, seules un nombre fini d'indéterminées jouant alors un rôle, et on va appliquer $0_{\overline{\underline{\mathrm{IV}}}}$ 15.1.16 $c)\Rightarrow b)$ : il faut vérifier que dans $k[X_n^\alpha]$, les $(X_n^\alpha)^p-P_n^\alpha(X_m^\beta,m<n)$ forment une suite régulière, et cela résulte d'un calcul de dimension trivial -

\hfill CQFD

La propriété que je n'ai pas su bien dégager est que toute extension algébrique $K/k$ est du type $K=k[T_\alpha]/(P_i)$, de sorte que si on relève les $P_i$ à un anneau artinien de corps résiduel $k$, soit $A$, $A[T_\alpha]/(\widetilde P_i)$ soit plat sur $A$.

$0_{\overline{\underline{\mathrm{IV}}}}$\quad\deleted{16.4.}\inserted{15}

Il manque des résultats sur les suites régulières, notamment le suivant auquel Quillen faisait une référence fantôme dans sa lettre sur $\Omega_{B/A}$.

% slanted
\margin{\unsure{regarder|Vs garder} EGA $\overline{\underline{\mathrm{IV}}}$ 19, à compléter éventuellement.}

\uline{prop} : \uline{Soit $A$ noethérien, $I$ un idéal}. \uline{Les conditions suivantes sont équivalentes}.

\[
\begin{aligned}
&\text{(i)}\quad\operatorname{Sym}_{A/I}(I/I^2)\xrightarrow{\sim}\operatorname{gr}_I(A)
&&\text{et }\ I/I^2\ \text{\uline{projectif sur} } A/I\\
&\text{(ii)}\quad \bigwedge I/I^2\xrightarrow{\sim}\operatorname{Tor}(A/I,A/I)
&&\text{et}\quad\text{''}\qquad\text{''}\quad\text{''}\quad\text{''}
\end{aligned}
\]

\uline{Si de plus $A$ est local}, et si $f_1..f_n$ est un \uline{système minimal de générateurs de $I$} (i é $n=\dim_k I\otimes k$), \uline{cela équivaut encore à}

\margin{\unsure{à|i}}

(iii) $f_1..f_n$ est \uline{une suite régulière}

(iv) le \uline{complexe de l'algèbre extérieure} $K(f,A)$ est \uline{une résolution} de $A/I$\inserted{\deleted{\unsure{et le reste par} \illegible{1}}} (?)

(v) $\operatorname{Ext}^i(A/I,A)=0$ si $i<n$

\uline{Démonstration} : (i) et (ii) étant de nature locale, on peut supposer $A$ local ; on peut aussi supposer $I\ne A$.

Il est clair que

\diagram{\begin{tikzcd}[ampersand replacement=\&,column sep=small,row sep=1pt]
\& \& \& \mathrm{(ii)}\\
\mathrm{(i)}\arrow[r,Leftrightarrow]\&\mathrm{(iii)}\arrow[r,Rightarrow]\&\mathrm{(iv)}\arrow[ur,Rightarrow]\arrow[dr,Rightarrow]\&\\
\& \& \& \mathrm{(v)}
\end{tikzcd}}

(car $\overline{\underline{\mathrm{IV}}}\Rightarrow$ que $\operatorname{Tor}^1(A/I,A/I)=I/I^2$ est $A/I$-libre)

--- l'implication \uline{$\mathrm{(ii)}\Rightarrow\mathrm{(iv)}$} : Soit $P_*$ une résolution projective de $A/I$ ; il existe une seule classe d'homotopie d'applications $K(f,A)\longrightarrow P_*$ (2) induisant l'identité sur $A/I$, et l'application $\bigwedge I/I^2=H_{\ast}^{\deleted{\unsure{*}}}(K(f,A)\otimes A/I)\longrightarrow\operatorname{Tor}(A/I,A/I)=H_*(\deleted{\unsure{K}}P_*\otimes A/I)$ déduite coïncide avec celle considérée en (ii), qui provient de la structure multiplicative de $\operatorname{Tor}(\ )$. [facile] - Soit $E$ le mapping cylinder de (2). (ii) signifie que $E\otimes A/I$ est acyclique, ce qui implique \unsure{pr|par} Nakayama que $E$ l'est, donc que (iv) est vrai.

--- l'implication \uline{$\mathrm{(v)}\Rightarrow\mathrm{(iii)}$} va être modulée.

\uline{prop} : \uline{Si} $I$ est un idéal de l'\uline{anneau local $A$}, $f_1..f_n$ un système minimal de générateurs de $\deleted{\unsure{A}}I$, $M$ un $A$-module de type fini et si

\[
\begin{aligned}
&\underline{\operatorname{Ext}^i(A/I,M)=0\quad\text{pour }i<n}\qquad\text{\uline{alors}}\\
&\underline{f_1..f_n\text{ est une suite }M\text{-régulière}}
\end{aligned}
\]
% archive sheet number: 62
